% SPDX-License-Identifier: Apache-2.0
\section{A Multiply-Accumulate}
\label{sec:x-mul}

\code{mul::<M>} is a product with its width stated. Rust cannot write
a width that is the sum of two others on stable, so the product does
not infer its width from its operands as the target languages do; it
takes it as \code{concat} does. The lowering extends both operands
to that width and multiplies there, which in Verilog is a
\code{*} whose result is the operands' width and in VHDL a
\code{resize} of a product twice as wide, so the two agree with the
Rust value, which wraps. The unit sums the products of the pairs it
is given, a cycle behind, and clears on request; it is checked under
nvc and Verilator against its trace, as Section~\ref{sec:x-checked}
describes. The Vreteno core does not use it: its multiply is a
sequencer of shifts and adds, thirty-two steps, which is a choice
about hardware rather than about the language.

\begin{figure}[htbp]
\centering
\input{mul_timing.tex}
\caption{The multiply-accumulate: a pair per cycle, the sum a cycle
behind, and a clear at either end.}
\label{fig:x-mul}
\end{figure}

\lstinputlisting[caption={\code{ex\_mul.rs}. Run with \code{bazel run //lib/examples:ex\_mul}.},label={lst:x-mul}]{lib/examples/ex_mul.rs}

\lstinputlisting[style=ascii,caption={What \code{ex\_mul} printed when the build ran it: the pairs and the sums, then the lowering.},label={lst:x-mul-out}]{out_mul.txt}
